
Dalam praktiknya, lazim pula melakukan pemeriksaan secara mental untuk menilai
apakah ada alasan untuk meyakini populasi yang mendasari
memiliki kemencengan sedang (jika $n < 30$)
atau memiliki pencilan yang sangat ekstrem ($n \geq 30$)
di luar apa yang kita amati dalam data.
Sebagai contoh, perhatikan jumlah pengikut
pada setiap akun individual di Twitter,
lalu bayangkan distribusinya.
Sebagian besar akun telah mengumpulkan
paling banyak beberapa ribu pengikut,
sementara sebagian yang relatif sangat kecil telah menghimpun
puluhan juta pengikut,
yang berarti distribusinya sangat menceng.
Ketika kita mengetahui bahwa data berasal dari distribusi
yang kemencengannya seekstrem itu,
diperlukan upaya tertentu untuk menentukan seberapa besar
ukuran sampel yang diperlukan agar syarat normalitas
terpenuhi.

%if we were sampling accounts from Twitter
%and examining the distribution of followers on the sampled
%accounts, we can expect that the vast majority of accounts
%will have fewer than 1,000 followers and that there
%will be some very extreme outliers who have tens of millions
%of followers.

%Distribution of the number of subscribers for
%      anyone who has uploaded a video to YouTube.
%      Most such individuals will have built little to
%      no following, while others will have amassed tens
%      of millions of subscribers.
%Generally, we do not presume you to always know when the
%underlying population has particularly extreme outliers.
%That~is, besides looking at the data itself,
%considering the mental check for whether particularly extreme
%outliers are likely to be a sanity check, not a formal check.

%\begin{figure}[h]
%  \centering
%  \Figure{0.8}{outliers_and_ss_condition}
%  \caption{Sample observations for
%      Example~\ref{outliers_and_ss_condition_ex}.}
%  \label{outliers_and_ss_condition}
%\end{figure}

%A more thorough sample size condition assessment would
%also consider two additional aspects beyond the core
%guidance above.

%The most nuanced checks are then when the sample size
%is very small -- and we have almost no observational data
%to allow us to check the condition.
%\begin{description}
%\item[Population knowledge.]
%    If we have information about the population beyond
%    what we've observed in the sample, we would consider
%    this information as well.
%    For example, if the sample size is under 30
%    and the population is known to be moderately skewed
%    (something difficult to detect with a small sample
%    in the observed data),
%    we might still not consider the 
%    For example, if the sample size is under 30 but the
%    population is known to be moderately skewed,
%    then the sample size condition is not reasonable.
%    Likewise, if the population is known to have particularly
%    extreme observations (examples below), then we may
%    require a particularly large sample size if we want
%    to use the normal model for $\bar{x}$.
%\item[Relaxing the extreme outlier condition.]
%    When the sample size gets very large,
%    we may even be able to overcome issues with
%    particularly extreme outliers.
%    However, there isn't clear guidance, and instead,
%    custom simulations can be helpful but are beyond
%    the scope of this book.
%\end{description}
%In this first course in statistics,
%you won't (and aren't expected to) have perfect judgement
%on when the sample size condition is or is not met.
%However, you are expected to be able to handle the
%clear cut cases based on the core guidelines.

%For those wanting to do more rigorous checks
%or for the situation that the ,
%then we add a 
%below are slightly more careful checks, it's convenient
%to break down 
%\begin{description}
%\item[Sample size under 30.]
%    If the sample size is less than 30, then we simply follow
%    the rule of thumb and there isn't
%    extreme skew in the data (usually punctuated by
%    extreme outliers), then we can proceed.
%\item[Sample size at least 30.]
%    If the sample size is at least 30 and there isn't
%    extreme skew in the data (usually punctuated by
%    extreme outliers), then we can proceed.
%\end{description}
%then it's generally reasonable to consider $\bar{x}$
%as following a nearly normal distribution.

%\Comment{Check the ``99\%'' and ``hundreds'' claim in the
%  income example below.}
%
%\begin{examplewrap}
%\begin{nexample}{Describe a couple populations that you know
%    would have particularly extreme outliers.}
%  Wealth distributions in many countries have
%      particularly extreme outliers.
%      For example, over 99\% of the population
%      has fewer than \$10 million saved,
%      while there are hundreds individuals in the
%      United States with over \$1~billion and who
%      are unlikely to be captured in even a moderate-sized
%      sample.
%
%  Distribution of the number of subscribers for
%      anyone who has uploaded a video to YouTube.
%      Most such individuals will have built little to
%      no following, while others will have amassed tens
%      of millions of subscribers.
%
%%  So while we won't be quizzing you on a variety of applications
%%  in this book, when you apply these skills elsewhere it is
%%  important to keep this consideration in mind and do some
%%  research if you aren't sure about outliers.
%\end{nexample}
%\end{examplewrap}
%
%Generally, we do not presume you to always know when the
%underlying population has particularly extreme outliers.
%That~is, besides looking at the data itself,
%considering the mental check for whether particularly extreme
%outliers are likely to be a sanity check, not a formal check.

%\begin{examplewrap}
%\begin{nexample}{Suppose we randomly sampled 20 individuals
%    from the United States and considered their incomes. %,
%    % which are shown in the following distribution:
%    However, the population is known to have particularly
%    extreme outliers, e.g. some individuals with incomes
%    above \$10 million.
%    No matter what we observe in the original 20 observations,
%    can you say whether we should proceed with modeling
%    $\bar{x}$ using a normal distribution?}
%  When we know the population distribution to have particularly
%  extreme outliers, then even if we observe no outliers in our
%  sample, we should not proceed to model $\bar{x}$ using
%  a normal distribution.
%
%  Generally, we do not presume you to always know when the
%  underlying population has particularly extreme outliers.
%  So while we won't be quizzing you on a variety of applications
%  in this book, when you apply these skills elsewhere it is
%  important to keep this consideration in mind and do some
%  research if you aren't sure about outliers.
%\end{nexample}
%\end{examplewrap}

%However, if one or more of clear outliers are present are evidently present,
%the guidelines around a reasonable minimum sample become murky.
%If 
%We'll see some other examples throughout the rest of this book,
%which will help in developing some intuition around this topic,
%but in many cases, data with .

\index{Teorema Limit Pusat!data normal|)}
